%% HMRS platform — Knowledge-Based Systems submission (elsarticle)
%% Hierarchical management framing: the system
%% manages documents at three levels (collection / document / segment) and
%% serves a tiered access path; retrieval algorithms are supporting components.
%% Compile: pdflatex main_20p && bibtex main_20p && pdflatex main_20p && pdflatex main_20p
%% AUTHOR ACTION (authors reserved these values; do not infer them):
%%   1. frontmatter below: author names/order, affiliation, city, country,
%%      corresponding-author e-mail, and the ORCID lines if any;
%%   2. n_back.tex: the two CRediT names, the funding wording and the
%%      acknowledgements;
%%   3. n_back.tex: confirm the generative-AI statement and figure provenance.
%% Default: double-column journal layout. main_20p_review.tex selects
%% a double-spaced layout from these same sources.
\ifdefined\HMRSReviewLayout
\documentclass[review,12pt]{elsarticle}
\else
\documentclass[final,5p,times,twocolumn]{elsarticle}
\usepackage{stfloats}
\fi

\usepackage[T1]{fontenc}
\usepackage[utf8]{inputenc}
\usepackage{booktabs}
\usepackage{array}
\usepackage{graphicx}
\usepackage{amsmath}
\usepackage{amssymb}
\usepackage{xcolor}
\usepackage{url}
\usepackage[colorlinks=true,linkcolor=blue,citecolor=blue,urlcolor=blue]{hyperref}
\usepackage{xurl}
\urlstyle{same}
% DOIs are emitted by the .bst as \path{...}; re-declare it so they
% typeset in the ambient roman face at 7.2 pt instead of 6.97 pt typewriter.
\DeclareUrlCommand\path{\urlstyle{same}}
\usepackage{caption}
\captionsetup{font=small,labelfont=bf}
\usepackage{algorithm}
\usepackage{algpseudocode}
\newcommand{\yes}{\checkmark}
\newcommand{\no}{--}

\graphicspath{{figs/}}
\setlength{\emergencystretch}{1.5em}

%% Allow figures and tables to stay near their first discussion.
\renewcommand{\topfraction}{0.90}
\renewcommand{\bottomfraction}{0.85}
\renewcommand{\textfraction}{0.08}
\renewcommand{\floatpagefraction}{0.75}
\renewcommand{\dbltopfraction}{0.90}
\renewcommand{\dblfloatpagefraction}{0.75}
\setcounter{topnumber}{3}
\setcounter{bottomnumber}{2}
\setcounter{totalnumber}{4}
\setcounter{dbltopnumber}{3}

%% Spacing for the double-column layout.
\ifdefined\HMRSReviewLayout\else
\setlength{\floatsep}{7pt plus 2pt minus 2pt}
\setlength{\textfloatsep}{9pt plus 2pt minus 2pt}
\setlength{\intextsep}{7pt plus 2pt minus 2pt}
\setlength{\dblfloatsep}{7pt plus 2pt minus 2pt}
\setlength{\dbltextfloatsep}{9pt plus 2pt minus 2pt}
\fi

%% Ragged bottoms: with nine floats the flushbottom setting stretches
%% underfull columns into heading-sized gaps (visible under the Threats and
%% Conclusion headings). Ragged bottoms leave the slack at the column foot
%% instead, which is the usual journal appearance for float-heavy pages.
\raggedbottom

\journal{Knowledge-Based Systems}

\begin{document}

\begin{frontmatter}

\title{HMRS: Hierarchical document management for evidence-traced retrieval}

\author[inst1]{Author Name\corref{cor1}}
\ead{author@example.com}
\author[inst1]{Co-author Name}
\cortext[cor1]{Corresponding author.}
\affiliation[inst1]{organization={Affiliation Department, University,
City, Country}, city={City}, country={Country}}

\begin{abstract}
Retrieval over evolving document collections depends on whether documents
are listed in the right collection and whether search can reach their
current content. The Hierarchical Management Retrieval System (HMRS) links
collection, document and segment state to catalogue, shelf, document and
segment views. A storage audit found at least one stored part for each of
100 target papers. A separate lifecycle check recovered a renamed document
using a cache-bypassing query.

In a scripted comparison on 72 questions across ten domains,
full route cards selected the labelled collection on 72 questions, versus
49 with nominally name-only cards. This contrast does not isolate the effect
of descriptions. The document-plus-chunk configuration returned the annotated
string for 48 questions, versus 36 for the reference-mode baseline and 41
for a configuration that passed longer document windows directly to the
judge. The paired comparison with the reference-mode baseline yielded 12
additional hits and no losses (two-sided exact test, $p=0.00049$, unadjusted).
Mean retrieval tool time was 15.0~s for the document-plus-chunk configuration
and 16.6~s for the reference-mode baseline, excluding answer generation.
These results distinguish collection selection from evidence access. They do not establish semantic support or answer accuracy: an exact
match can occur in the wrong document or context.
\end{abstract}

\begin{keyword}
Knowledge document management \sep Document classification \sep
Hierarchical organisation \sep Tiered retrieval \sep Evidence verification \sep
Model Context Protocol
\end{keyword}

\end{frontmatter}

%% ============================================================ body (v2 order)
\section{Introduction}
\label{sec:intro}

Enterprise documents preserve regulations, failure analyses and operating
procedures. To answer a question from these materials, a retrieval system
must reach the right collection, identify a document and expose the passage
that supports the answer. Each step depends on earlier management decisions.
A collection description guides the search, a document must be listed on
the selected shelf, and the index must point to its current location.
An accurate passage ranker has little opportunity to recover evidence that
the access path never reaches.

Document acquisition, organisation and retrieval are often evaluated at
different boundaries. Enterprise content management studies document
governance \cite{paivarinta-ecm-2005,tyrvaainen-ecm-2006}, while metadata
quality and tagging determine how users discover those documents
\cite{park-metadata-quality-2009,golder-tagging-2006}. Distributed
information retrieval addresses the selection of collections
\cite{callan1995searching,si2003redde}. Retrieval-augmented generation
(RAG) and agentic retrieval add passage ranking and tool-guided evidence
access \cite{rag-lewis,dpr,modular-rag,mcp}.
HMRS connects these responsibilities. It makes collection membership,
document metadata and source references available from ingestion through
query time, so that an operator can trace a failed search to the stages
that supplied its candidates.

The Hierarchical Management Retrieval System (HMRS) is a self-hosted
knowledge document platform. An ingest agent uses parsed content to select
a collection and describe the document. Management operations update the
files, tree registry and library metadata, and reconciliation checks them
against one another. At query time, a catalogue-first search uses collection
and document descriptions to reach source passages. A pinned search begins
within a selected collection, and an optional server-side evidence gate
returns scored passages with source coordinates. The access policies use the
same managed state, so their choices can be traced independently of parsing
or index maintenance.

\noindent\textbf{Contributions.}
\begin{enumerate}
\item \textbf{A management model for the state behind retrieval.} HMRS
organises knowledge documents as collections, documents and segments, and
connects these managed objects to catalogue, shelf, document and segment
views. Collection membership, descriptions and source references link the
management of documents to the choices made during evidence access
(Section~\ref{sec:system}).
\item \textbf{A content-informed document lifecycle.} The acquisition
workflow uses parsed content to select a collection and derive document
descriptions and tags. Management operations and reconciliation checks
connect the document files, tree registry and per-library metadata, while
index and retrieval probes help expose stale state after changes
(Sections~\ref{sec:acquisition}--\ref{sec:organization}).
\item \textbf{Agent access across management levels.} MCP interfaces expose
the catalogue, document descriptions, source documents and indexed passages
to an agent. Catalogue-first and pinned policies let the agent narrow a
search using managed descriptions or start within selected collections
(Section~\ref{sec:retrieval}).
\item \textbf{A source-traceable evidence gate.} An optional server-side
gate fetches candidate text, scores its passages against the question, and
returns admitted passages with source references and error information.
Candidate ordering is a separate decision made by the calling workflow or
hybrid tool. Answer instructions distinguish a completed search without
support from an incomplete execution
(Sections~\ref{sec:gate} and~\ref{sec:lanes}).
\end{enumerate}

\subsection{Research questions and reading guide}
\label{sec:rq}

\textbf{RQ1} examines document acquisition, collection membership and
lifecycle consistency. \textbf{RQ2} measures collection selection and
document visibility under catalogue-first access. \textbf{RQ3} compares
target-string containment and retrieval cost across the tested evidence
configurations. \textbf{RQ4} compares the order of retained candidates and
examines the admission scores available in the archive.
\textbf{RQ5} examines abstention, runtime, tool cost and execution failures
at the answer layer.

Section~\ref{sec:related} positions HMRS within related work.
Sections~\ref{sec:system}--\ref{sec:implementation} describe the management
model and its implementation. Section~\ref{sec:eval} reports the evaluation,
followed by the discussion, limitations and conclusions in
Sections~\ref{sec:discussion}--\ref{sec:conclusion}.

\section{Related work}
\label{sec:related}

\subsection{Knowledge management and document acquisition}

Recent work on knowledge-based systems combines language models with curated
knowledge \citep{yang-kbs-llm-2025}. DBpedia publishes structured
information from Wikipedia as linked data \citep{auer-dbpedia-kb-2007}.
Enterprise content management governs the documents through which
organisations retain and share knowledge
\citep{paivarinta-ecm-2005,tyrvaainen-ecm-2006}. Knowledge creation theory
also examines how organisations articulate tacit knowledge
\citep{nonaka-seci-1994}.
Metadata, tagging and human review affect which materials users can find and
trust \citep{park-metadata-quality-2009,golder-tagging-2006,mosqueira-rey-hitl-2022}.
Text categorisation supports placement \citep{sebastiani2002machine}.
Layout-aware parsing and MinerU recover structured content for indexing
\citep{zhong-publaynet-layout-2019,huang-layoutlmv3-docai-2022,mineru}, while
near-duplicate detection and provenance help maintain the collection
\citep{manku-near-dup-2007,missier-prov-2013,cheney-provenance-2009}.
HMRS connects these operations by making the results of acquisition and
placement visible to the workflow that selects evidence.

Gallofr\'e Oca\~na and Opdahl \citep{ocaa2023software} describe how to
integrate AI and knowledge bases in a newsroom platform. Related systems
ground dialogue in a graph, build open-domain question answering for
low-resource languages, or support retrieval by bank supervisors
\citep{yang2024building,budur2024building,aarab2024system}. LLMs can also
populate domain ontologies \citep{ciatto2025large}. HMRS applies this systems
perspective to changing document collections. Its focus is the managed
state a retrieval workflow reads and an operator can inspect after a
lifecycle operation.

\subsection{Retrieval, organisation and agentic search}

\begin{table*}[!t]
\centering
\small
\setlength{\tabcolsep}{3.5pt}
\renewcommand{\arraystretch}{1.02}
\caption{Representative systems compared along the dimensions examined here.
A dash indicates a dimension outside the cited work's main contribution.}
\label{tab:related-comparison}
\begin{tabular}{@{}>{\raggedright\arraybackslash}p{0.13\textwidth}>{\raggedright\arraybackslash}p{0.245\textwidth}>{\raggedright\arraybackslash}p{0.21\textwidth}>{\raggedright\arraybackslash}p{0.215\textwidth}>{\raggedright\arraybackslash}p{0.13\textwidth}@{}}
\toprule
System & Hierarchy and scope & Construction and lifecycle & Evidence handling & Access surface\\
\midrule
RAPTOR \citep{raptor} & cluster tree over document text & built offline at index time & --- & tree retrieval; prototype\\
GraphRAG \citep{graphrag} & entity--relation graph, community summaries & pre-generated offline & --- & graph retrieval; prototype\\
LightRAG \citep{lightrag} & low- and high-level graph retrieval & incremental graph updates & --- & graph retrieval; prototype\\
DeepRead \citep{deepread} & per-document structured Markdown, section coordinates & structure built at ingestion & agent locates then reads sections & agentic search system\\
EPRAG \citep{kinalioglu-kbs-eprag-2026} & multi-source enterprise corpus & --- & evidence-state diagnosis; abstention or clarification & enterprise retrieval\\
HMRS (this work) & managed collection--document--segment state & content-based placement; lifecycle operations; three-source reconciliation & optional gate, separate from ordering & self-hosted; MCP layer\\
\bottomrule
\end{tabular}
\end{table*}

RAG uses retrieved passages to inform generation \citep{rag-lewis}. Dense
retrieval ranks passages in a learned representation space \citep{dpr}; BM25
uses lexical signals \citep{bm25}. Enterprise architectures combine semantic
retrieval with vector storage and LLMs \citep{ghali-kbs-retrieval-2025}.
Hybrid search combines ranking signals, for example with reciprocal-rank
fusion \citep{rrf-cormack}, and BGE-M3 provides multilingual and
multi-granularity embeddings \citep{bge-m3}. Structured retrieval organises
the units a ranker searches. GraphRAG builds entity graphs and community
summaries \citep{graphrag}. LightRAG searches an incrementally updated
graph at two retrieval levels \citep{lightrag}, and RAPTOR recursively
groups text into a cluster tree \citep{raptor}. RMR-SST represents multimodal engineering
documents as semantic trees \citep{li-kbs-rmrsst-2026}. Dense X Retrieval
indexes propositions \citep{chen2024dense}, while TEIE-RAG combines
hierarchical semantic chunking with source highlighting for engineering
information extraction \citep{shen-kbs-teierag-2026}.
Table~\ref{tab:related-comparison} locates HMRS among these designs. Its
hierarchy represents collection ownership and document structure, so a
query reads a maintained record of which collections and documents exist.
Index maintenance complements this hierarchy by keeping passages reachable
after document changes \citep{xu2023spfresh}.

Distributed IR selects collections from relevance estimates before searching
them \citep{callan1995searching,si2003redde}. HMRS applies the same
select-then-search principle to collection and document descriptions that
operators can maintain. Tool-using agents can choose and iteratively refine
what they retrieve \citep{react-yao,toolformer}; MCP provides standard
interfaces to these tools \citep{mcp}. Modular RAG studies how these
workflows combine retrieval components \citep{modular-rag}. DeepRead
illustrates structure-aware agent access by indexing paragraphs with section
coordinates and reading selected sections \citep{deepread}. HMRS gives
an agent a catalogue and shelf to determine which documents to open, then
exposes document and passage access through the same tool layer.

In engineering design, structured facts extracted from patents support RAG
\citep{siddharth2024retrieval}. Other approaches change which evidence a
query can reach through soft-prompt tuning, semantic query expansion and
stateful graph exploration
\citep{peng2024soft,massai2023evaluatio,tao2025fine}. HMRS exposes the document
state on which collection and passage selection depend.

\subsection{Verification, correction and abstention}

MemGPT and A-MEM maintain agent interaction records and notes
\citep{memgpt,amem}. HMRS maintains a corpus of source
documents and the descriptions, ownership and indices used to find them.
Knowledge organisation standards such as SKOS \citep{skos} and Dublin Core
\citep{dcmi-terms}
give conventions for vocabularies and descriptive metadata. At ingestion,
the HMRS workflow derives descriptions and tags from document content.

Self-RAG and CRAG assess evidence or correct retrieval during generation
\citep{selfrag,crag}; FLARE triggers retrieval as more evidence is needed
\citep{flare}. Work on selective prediction under domain shift and
uncertainty reporting examines when to withhold an answer
\citep{kamath,feng-abstain,yin-knowwhat}. Citation studies examine whether
references support the statements in the answer
\citep{liu-verifiability,gao-citations}.
EPRAG diagnoses evidence states to choose among response, clarification and
abstention in enterprise retrieval \citep{kinalioglu-kbs-eprag-2026}.
HMRS makes verification a server-side operation. The gate fetches
candidate text, scores passages and returns surviving text with provenance,
so an agent can inspect the source material behind the result.

Judge studies report preference bias, variable agreement and ranking
instability \citep{zheng-llmjudge,sun-urara,panickssery-selfpref}.
Logic-grounded metamorphic testing checks whether model decisions remain
consistent under equivalent input transformations \citep{zhou2026lgmt}.
HMRS tests evidence admission, candidate order and answer-stage reporting
as distinct decisions.

\subsection{Research gap}
\label{sec:gap}

RMR-SST and TEIE-RAG structure engineering document content for retrieval
and evidence access \citep{li-kbs-rmrsst-2026,shen-kbs-teierag-2026}. EPRAG
selects a response action once the candidate evidence is available
\citep{kinalioglu-kbs-eprag-2026}. HMRS studies the collection state that
supplies evidence to such workflows (Table~\ref{tab:related-comparison}
and Supplementary Table~S11). It records where documents belong, exposes
the descriptions used to select them, and lets operators check whether
the index reflects lifecycle changes. The evaluation follows this path from
post-ingest state and catalogue visibility through candidate selection and
verifier ordering.

\section{System overview and management model}
\label{sec:system}

HMRS manages the state a search reads. The parsing service recovers document
content, the web application manages the collection and its files, and the
MCP server exposes these operations to agents \cite{mcp}. A document's
placement determines the scope in which it is visible. Its description
guides a catalogue-first search, and its index entries make passages
discoverable to lexical and dense search. To check these connections, HMRS
compares document files, the tree registry and per-library metadata after
selected management operations.

\subsection{Management levels and the tiered access path}
\label{sec:levels}

The management model has three levels: collections carry ownership and scope,
documents carry content and descriptive metadata, and segments carry the
text units used for passage retrieval or scoring. Four access views expose
those levels. The catalogue lists collections, the shelf lists documents
within a collection, the document view opens a body and the segment view
exposes passages and their source references. Table~\ref{tab:levels} links each
managed level to the operations and checks that exercise it.

\begin{table*}[t]
\centering
\small
\setlength{\tabcolsep}{4pt}
\renewcommand{\arraystretch}{1.15}
\caption{Management levels and the tiered access path. Each level names the
objects it maintains, the operations that write them, the state that describes
them, the tier that reads them, and the measurement that exercises it
(evidence labels E1--E9 denote the evaluation campaigns of
Section~\ref{sec:eval}; Table~S17 records their design). The selected access
policy determines which tiers are visited; a pinned
search can bypass catalogue navigation.}
\label{tab:levels}
\begin{tabular}{@{}>{\raggedright\arraybackslash}p{0.12\textwidth}>{\raggedright\arraybackslash}p{0.21\textwidth}>{\raggedright\arraybackslash}p{0.18\textwidth}>{\raggedright\arraybackslash}p{0.12\textwidth}>{\raggedright\arraybackslash}p{0.12\textwidth}>{\raggedright\arraybackslash}p{0.13\textwidth}@{}}
\toprule
Level & Objects and operations & State description & Read tier & Check & Evidence\\
\midrule
Collection & libraries and sub-libraries; agent assignment from parsed content; description and tag derivation; catalogue listing & library registry, per-library metadata, library-level index & catalogue & catalogue inventory and description round-trip & E1, E2, E3\\
Document & documents and their parts; duplicate review; move, rename, delete; forced re-index & document files, tree registry, per-library metadata & shelf & reconciliation on selected operations & E2\\
Segment & structural split; dense and lexical indices; optional evidence gate & chunk identifiers, per-segment scores with provenance & document and segment & index coverage and post-index probe & E1, E5--E9\\
\bottomrule
\end{tabular}
\end{table*}

The ingest agent uses catalogue descriptions to place parsed material and
derives document descriptions and tags from its content. Sub-libraries are
tree nodes with their own descriptions and indices. At the document level,
duplicate review identifies candidates and a lifecycle operation can move,
rename or delete a document. At the segment level, search returns indexed
chunks, while the optional evidence gate fetches candidate text and scores
query-conditioned passages with provenance.

Figure~\ref{fig:levels} uses the library analogy to show the transitions
between catalogue, shelf, document and segment views. A pinned search can
start inside selected collections and access indexed chunks directly.
Figure~\ref{fig:arch} shows how both paths reach the same document and
index services. The catalogue and shelf can be inspected independently:
an absent collection and an omitted document appear at different stages
of the access trace.

\begin{figure}[t]
\centering
\includegraphics[width=\columnwidth]{fig_library_analogy_vector}
\caption{Library analogy for a catalogue-first HMRS access path. The library catalogue, shelf, book and passage correspond to the collection catalogue, document listing, document body and segment. Evidence scoring is optional; pinned search may bypass the catalogue and shelf.}
\label{fig:levels}
\end{figure}

\setcounter{dbltopnumber}{1}
\begin{figure*}[t]
\centering
\includegraphics[width=\textwidth]{fig_service_architecture_vector}
\caption{HMRS service architecture. The browser and agent/CLI access the web and MCP interfaces. Parsing, management and retrieval operate on shared document files, the tree registry, library metadata and search indices. Laya scoring is optional, and the Neo4j graph supports navigation.}
\label{fig:arch}
\end{figure*}

Figure~\ref{fig:workflow} follows a document from acquisition and
organisation to evidence-based reporting. The ingest workflow governs
content review and placement. The query workflow determines whether to
consult the catalogue or search a pinned collection, and whether to use
Laya to score passages. The answer instructions distinguish a completed
review's searched scope and blind spots from an incomplete review caused
by tool failure.
Hybrid retrieval is also available as a tool and can be composed from
ordinary search calls. Experience retrieval and persona distillation reuse
the document infrastructure and define their own retrieval procedures.

Table~\ref{tab:functions} maps acquisition, organisation, lifecycle, query
and experience operations to the interface or workflow that checks them.
This allocation locates decisions at the appropriate boundary. For example,
the ingest agent chooses a collection, the save service validates the
destination identifier, and an audit checks the stored result.

\begin{table*}[t]
\centering
\small
\setlength{\tabcolsep}{4pt}
\caption{Management functions, the enforcement or workflow check associated
with each, and the evaluation that exercises it. The table separates API/MCP
controls from agent procedures and campaign audits; semantic choices such as
which collection owns a document are made and recorded by the workflow.
Evidence labels E1--E9 denote the campaigns in Section~\ref{sec:eval}.}
\label{tab:functions}
\begin{tabular}{@{}>{\raggedright\arraybackslash}p{0.13\textwidth}>{\raggedright\arraybackslash}p{0.37\textwidth}>{\raggedright\arraybackslash}p{0.30\textwidth}>{\raggedright\arraybackslash}p{0.11\textwidth}@{}}
\toprule
Group & Operations & Platform or workflow check & Evidence\\
\midrule
Acquisition & submit, parse (MinerU), quality review, duplicate review, structural split, store, index & parse-task state; workflow quality checks; entry-point-specific deduplication and indexing; post-index probe & E1--E2, S9\\
Organisation & create library or sub-library, content-informed placement, descriptions and tags, duplicate review & API validates identifiers and required fields; workflow checks descriptions and tags; audit compares files, tree and library metadata & E1--E2\\
Lifecycle & move, rename, delete, force re-index, graph build & operation-specific API updates; reconciliation audit detects selected stale or inconsistent state & E2\\
Query & tier selection, vector, keyword, catalogue and two-stage search, gate invocation & endpoint enforces tool allow/deny lists, timeouts and turn limits; invoked gate returns survivor provenance & E3, E5--E9\\
Experience & extract, approve, tier, search & experience tools apply their own retrieval, verification and credibility filters & E4\\
\bottomrule
\end{tabular}
\end{table*}

\begin{figure*}[t]
\centering
\includegraphics[width=\textwidth]{fig_workflow_vector}
\caption{Illustrative HMRS workflow from document acquisition to evidence-based reporting. Documents are parsed when needed, organised into collections and indexed according to the entry point. A question can use catalogue navigation or pinned search; optional Laya scoring precedes evidence inspection. A completed review without support is distinct from an incomplete search or tool failure, which may occur at any stage.}
\label{fig:workflow}
\end{figure*}

\subsection{Data model and object lifecycle}
\label{sec:datamodel}
\label{sec:lifecycle}

Documents are stored as Markdown with associated assets. A managed
collection's state spans the document files, \texttt{.tree-fs.json} and the
library's \texttt{.knowledge-base.yml}. The ingest workflow distinguishes
parsed, stored and indexed content and records completion after a retrieval
probe. The MCP save path can start background indexing; the web parsed-file
save path exposes indexing as a separate operation. Reconciliation compares
the three representations and identifies selected stale-index conditions.
A move, rename or re-index can therefore be checked through the resulting
file path, metadata and retrieval response.

\section{Acquisition and automatic classification: from parsed document to
owned collection}
\label{sec:acquisition}

Acquisition establishes what a later search can find and where it can
look for it. HMRS parses PDFs, scanned images, Word and Excel files through
an asynchronous MinerU service \citep{mineru}. The MCP parse tools return
pollable task identifiers. Markdown and plain text take the direct save
path. For parsed inputs, the ingest agent reviews the recovered body before
selecting a collection by topic.

The ingestion workflow connects parsing, review, optional splitting,
placement, storage and indexing. The source content informs the placement
decision, and the destination determines the scope of subsequent access.
The workflow records the applicable artefacts: a parse task identifier,
a quality report, a content fingerprint, a split manifest, a storage
record and the destination library. A content-based retrieval probe checks
that indexing exposes the document to later queries. This check follows the
save acknowledgement and tests a different transition.

\subsection{Quality gates and ingest lifecycle}

The upload service checks sibling files for exact duplicates using content
hashes. A separate near-duplicate tool returns pairs above a similarity
threshold for review. The ingest agent checks recovered text for OCR noise,
binary residue, missing section bodies, whitespace collapse and language
mismatch. It selects a library whose scope matches the recovered subject and
derives descriptions and tags from the same body. The owning collection and
its listing can then be checked after storage.

On oversized documents, the optional structural splitter groups recognised
units such as headings, Markdown tables, lists, block quotes and fenced code.
The research-paper campaign used a 30,000-character limit; later configurations used
60,000. The splitter aims to fill 85\% of the configured window and records
oversized indivisible units where applicable. Structure recognition covers
Markdown units, with formula-heavy pages depending on the quality of the
parsed representation.

The web parsed-file route persists Markdown and images. The MCP save tool
can also schedule vector and lexical indexing in the background. Graph
construction is a separate operation. Post-index checks can compare
expected identifiers with the recorded chunks and probe retrieval for
content from the saved document.

\subsection{Acquisition evidence and accounting limits}

The research-paper campaign provides an acquisition audit across five
libraries. Its target, indexed parts and observed retrieval coverage use
different units and are reported with their dates in Section~\ref{sec:e1}.

In the separate multilingual ingestion test, a script sent 400 Wikipedia
pages per language to preassigned libraries. Their listings contained
87.5\% of expected English titles and all expected Chinese and Japanese
titles. This test checks listing coverage after a known placement decision.
Section~\ref{sec:e1} describes the listing and self-retrieval checks on the
same libraries.

\section{Collection and document management: organisation, classification and
lifecycle}
\label{sec:organization}

Collection membership defines the scope available to a query. HMRS uses
the subject of the content for placement, preserves the source assignment
across split parts, and exposes descriptions before an agent opens a
document body. The collection and document operations in
Table~\ref{tab:levels} maintain these choices over the life of the corpus.

\subsection{Classification, metadata, tags and duplicate review}
\label{sec:org-placement}
\label{sec:org-registries}
\label{sec:org-tags}
\label{sec:org-dedup}

The catalogue returns collections and sub-libraries with descriptions.
Each shelf shows document titles and descriptions within its selected scope.
The ingest agent uses that scope and the parsed content to select a library,
then writes descriptions and tags grounded in the body. Duplicate review
presents candidate pairs for inspection because adjacent parts may share
metadata while retaining distinct source sections.

HMRS supports library and sub-library creation, metadata edits, moves,
renames, deletion and forced re-indexing. A move updates the document path
and registry entries, while re-indexing updates the searchable representation.
Integrity checks compare document existence, path and index metadata across
the files on disk, the \texttt{.tree-fs.json} registry and the library's
\texttt{.knowledge-base.yml}. These checks give operators observable
outcomes for the state changes made by each operation.

Sub-libraries keep their descriptions and indices at their own tree nodes.
When the caller supplies a parent scope, the vector service resolves the
parent and its descendants into the collections searched. The lexical stage
of the two-stage service likewise resolves descendant identifiers. The
tested flat union calls explicit library scopes, which fixes the collections
compared in the evaluation.

\subsection{Lifecycle operations, graph and operational knowledge}
\label{sec:org-lifecycle}
\label{sec:org-graph}
\label{sec:org-experience}

A rename trial exposed a lifecycle defect: queries found chunks under the
old document path until a forced re-index. The current rename paths
synchronise the old and new index entries and record failures separately
from the metadata update. A post-operation probe can then check whether
queries resolve to the new path.

The document graph links libraries, documents and tags for navigation and
candidate expansion. Operational lessons are also stored as retrievable
documents. Experience tools track reviews, use and freshness and use
credibility tiers to order results. Section~\ref{sec:eval} reports
operational-experience cases alongside the management experiments.

\subsection{Organisation contract}
\label{sec:org-integrity}
\label{sec:org-summary}

The organisation contract links each document to its current owning
collection, descriptive metadata and indexed path. The checks in
Section~\ref{sec:e2} test these relationships on selected lifecycle
operations by reconciling the three sources, reviewing document descriptions
and rerunning retrieval probes. The resulting traces locate inconsistencies
at the file, registry or index stage.

\section{The tiered access path: from collection to admitted segment}
\label{sec:retrieval}
\label{sec:lanes}

A query reaches evidence by narrowing its scope. Catalogue-first search
selects a collection, examines the descriptions on its shelf and opens
selected documents or passages. Pinned search starts within specified
collections and uses their indices to propose candidates. Both paths can
call the server-side Laya evidence gate. This design gives an agent
explicit choices at collection, document and passage granularity and
preserves the source identifiers needed to inspect those choices.

Figure~\ref{fig:tiered-algorithm} and Algorithm~\ref{alg:tiered} show how
an access policy selects, deduplicates and reads candidates, with or without
the evidence gate. The resulting pack is an input to answer assessment.

Access policies choose which operations to run. The intended workflow distinguishes
a completed review without support from an incomplete run with possible
partial evidence. This distinction depends on observable failure status;
the global-search defect in Section~\ref{sec:e7} shows it is not guaranteed
when a collection is silently skipped.

\begin{figure*}[t]
\centering
\includegraphics[width=\textwidth]{fig_tiered_drawio_20261116}
\caption{Tiered access workflow. Candidate references come from pinned
collections, catalogue navigation or a combination. The optional Laya gate
fetches and scores passages. A completed review without support reports
the searched scope. An incomplete operation reports failures and keeps
any available evidence.}
\label{fig:tiered-algorithm}
\end{figure*}

\begin{algorithm}[t]
\caption{Tiered access with optional evidence gate}
\label{alg:tiered}
\begin{algorithmic}[1]
\Require question $q$; catalogue $\mathcal{C}$; threshold $\tau$; gate flag $g$
\State $D \gets \textsc{GatherCandidates}(q,\mathcal{C})$ \Comment{pinned or catalogue-first}
\If{$g$}
  \State $S \gets \textsc{LayaJudge}(D,q,\tau)$ \Comment{fetch, score and return survivors}
\Else
  \State $S \gets \textsc{ReadSelected}(D)$
\EndIf
\If{execution incomplete}
  \State \Return incomplete status, searched scope, and available evidence
\ElsIf{completed review finds no support}
  \State \Return scoped not-found (searched scope, blind spots)
\Else
  \State \Return agent-synthesised answer with source references
\EndIf
\end{algorithmic}
\end{algorithm}
The policy instructs the agent to report its searched scope and remaining
blind spots when it reaches its call budget. The chat endpoint controls
timeouts and turn limits. Workflow instructions specify the retrieval
sequence and batch sizes.

\subsection{Tier policies and the evidence gate}
\label{sec:trace}
\label{sec:gate}
\label{sec:gate-primitive}
\label{sec:twostage}

The knowledge-base chat interface defines two policies. The search policy
starts from pinned collections and calls two-stage or vector search. The
librarian policy starts from the all-library catalogue, uses shelf
descriptions to choose document references and directs the agent to score
them in sequential batches before reading the survivors. The two-stage
tool uses lexical search and optional graph expansion to form a candidate
pool, applies dense ranking and returns candidates for inspection.

The MCP server also registers \path{kb_hybrid_search}, which merges vector
and catalogue candidates, optionally applies Laya, and sorts retained
documents by verifier score when judging is enabled. The librarian policy
directs the agent to use the catalogue and shelf tools directly. A separate
controlled hybrid experiment used per-library vector top-30 at cosine
threshold 0.35, catalogue candidates, fusion and a cap of twelve. The
ordering analysis in Section~\ref{sec:e6} replays historical hybrid pools
to compare ways of using the recorded scores.

The local Laya classifier \cite{laya}, exposed as \path{kb_laya_judge},
segments candidate text. Its \texttt{noul} output is a passage score on a
0--1 scale; passages at or above the threshold (0.5 by default) are eligible
for admission. It is not a calibrated probability that a passage supports
an answer claim. The gate groups survivors by source and position,
deduplicates them and applies a character cap. Where available, survivors
carry collection and document identifiers, file path, part or section, and
line range in the parsed text. These locators are not bound to immutable
versions. Fetch or scoring errors are reported separately from low scores.
Calling workflows may apply their own ordering. Table~\ref{tab:lanes}
summarises the policy choices.

\begin{table*}[t]
\centering
\small
\setlength{\tabcolsep}{5pt}
\caption{Access policies over the tiered path. Modes differ in how candidate collections and documents are selected.
The hybrid row describes the experimental composition used in the evaluation.}
\label{tab:lanes}
\begin{tabular}{@{}>{\raggedright\arraybackslash}p{0.11\textwidth}>{\raggedright\arraybackslash}p{0.19\textwidth}>{\raggedright\arraybackslash}p{0.24\textwidth}>{\raggedright\arraybackslash}p{0.16\textwidth}>{\raggedright\arraybackslash}p{0.20\textwidth}@{}}
\toprule
Policy & Tiers traversed & Retrieval affordance & Budget type & Intended terminal report\\
\midrule
Search & pinned collections, segment & Two-stage or vector search & Endpoint turn limit & Answer or scoped not-found\\
Librarian & catalogue, shelf, document, segment & Catalogue navigation and passage scoring & Workflow instructions & Answer or searched scope\\
Hybrid & catalogue and segment & Combined vector and catalogue candidates & Configuration-specific candidate cap & Candidates or judged survivors\\
\bottomrule
\end{tabular}
\end{table*}

Response transport also determines which candidates the agent can see.
The catalogue-truncation case exposed an incomplete preview before the agent
could inspect the whole shelf. The shelf tool now caps its response and
supplies grouped and paginated views, allowing the caller to request the
remaining entries. Section~\ref{sec:eval} reports the repair check.

\subsection{Honest abstention}
The answer workflow asks the agent to report absence after a completed
review that finds no support, including the scope it searched. A plausible
candidate may still lack the required evidence. If a tool is unavailable
or the search times out, the workflow reports the incomplete execution
and any usable evidence. Section~\ref{sec:eval} reports how the agents
handled these terminal conditions in the observed trials.

\section{Agent interface and implementation}
\label{sec:implementation}
\label{sec:impl}

\subsection{MCP tools and lane policy}
\label{sec:impl-tools}

MCP exposes the management groups in Table~\ref{tab:functions} to agent
harnesses. Long-running parsing, forced re-indexing and graph construction
return task identifiers that a caller can poll. Catalogue, document read,
search and passage-scoring operations expose different views of the
resulting state. An operator can verify a write by probing its resulting
collection and document paths, then re-index if needed. Figure~S1 shows the tool layer and
its management and retrieval services.

The knowledge-base chat endpoint mounts the MCP server and supplies
read-oriented tool sets. Administrative writes use management tools and
their own permissions. The endpoint sets time and turn limits, while the
policy instructions direct the retrieval sequence and set batch sizes.
A separate KB agent endpoint adapts multiple harnesses for management and
other knowledge-base tasks and supports asynchronous task execution.
Persona distillation uses the same storage infrastructure with its own
retrieval and synthesis path.

\subsection{Runtime and deployment}
\label{sec:impl-runtime}

The backend uses Python 3.12 and FastAPI. ChromaDB stores cosine-space HNSW
collections with normalised 1024-dimensional BGE-M3 embeddings
\cite{chroma,hnsw,bge-m3}. BM25 supports lexical search \cite{bm25}, Neo4j
stores the document graph \cite{neo4j} and MinerU provides parsing and OCR
\cite{mineru}. The local Laya worker queues scoring jobs to serialise GPU
execution.

The benchmark ran on one self-hosted Windows workstation with an RTX
5060 GPU (8 GB). A recovery script checks the API, web and MCP services
before probes. Service readiness does not establish that all optional
retrieval components are available, as the deployment checks
illustrate (Section~\ref{sec:sysarms}).

Use of the experimental corpora remains subject to their data-use
constraints (see Data availability).

%% Evaluation design and E1--E3. Section labels follow the historical source.
\section{Evaluation}
\label{sec:eval}
We evaluate what is stored and remains accessible after a change (E1--E2,
E4), what the retrieval path selects and returns (E3, E5, E8), and what
answer diagnostics and deployment assertions establish (E6--E7, E9).
Tables~S9--S10 and S17 give the protocols and denominators. The separate
experiment source map links the descriptive campaign names to their
archives. We do not pool engineering assertions with benchmark results.

\subsection{Experimental design and measures}
\label{sec:eval-setup}

\textbf{Corpora and questions.} The routing and evidence-pack benchmark
uses the same 72 in-domain questions across ten domains in three campaigns:
the Frozen aggregate (138 documents stored as 879 parts), the Follow-up
(906 parts), and the Main comparison (seven arms on a restarted stack).
The Main comparison is the basis for the primary contrasts; the other
campaigns show how the same labels vary across corpus and index states. Each complete arm also has
four out-of-domain probes (R1--R4); two cross-domain placeholders were not
scored. Developers authored the questions and target annotations. A separate
paper corpus has 100 target papers and ten question-specific source IDs.
The historical industrial-corpus ranking diagnostic has 16 questions on
42 documents, including twenty near-duplicate telemetry reports; Q16 shares
a token with its gold filename. Tables~S9--S10 and S17 document these and
the separate multilingual listing checks.

\textbf{Measures.} The route metric checks which collection the script
selected. Containment checks for the annotated literal string anywhere in
the returned pack, even in a non-gold document. The historical ranking
analysis uses document Hit@5; the answer diagnostics check literal source
IDs, annotated terms (at least $\lceil 3n/5\rceil$ of $n$) and marker syntax.
These measures do not substitute for judged relevance or claim support
\citep{craswell-msmarco-2021,jarvelin-cumulated-gain-2002,es-ragas-2024,saad-falcon-ares-2024}.
Wilson intervals are marginal summaries; questions are paired by ID for
two-sided exact McNemar tests. Tests are exploratory and unadjusted.
Each answer condition was run once; retries belong to the same condition,
not independent trials. The historical grids used Claude Code through the
platform chat API; the four-question engine contrast in Table~S16 is a
separate instructed case.

\subsection{E1: Acquisition and parsing}
\label{sec:e1}
The storage audit found at least one stored part for each
of the 100 target papers. The five libraries contained 328 parts, including
six outside the target set. All 328 passed the recorded description check.
Neither check measures retrieval coverage. A filename-stem diagnostic found
all target identifiers, but most searches were unscoped and chunk copies
could count. A separate ten-question vector diagnostic ranked each gold
document first within its preassigned library; it did not test automatic
routing or passage support. Table~S9 preserves the audit details.

In a separate batch fixture, the warmed parser processed three
PDFs in 38.7~s. Those three parsed documents and two directly ingested
Markdown documents were all retrieved at rank one by document-specific
probes. This checks both acquisition paths on five inputs, not parser
accuracy or general retrieval recall.

\subsection{E2: Lifecycle consistency}
\label{sec:e2}
The single-document rename fixture compared three observations:
the stored metadata and vector entries showed the new name,
but the repeated query initially showed the old one. A cache-bypassing query
returned the new name, and a probe after forced reindexing confirmed it.
This distinguishes cached output from index state in one fixture; it does not
establish concurrent-write or crash-recovery safety.
An older rename test had found duplicate chunks before forced reindexing;
the two checks describe different system states. Table~S9 and the historical
inventory retain the separate reconciliation, sub-library, tagging and
graph checks.

\subsection{E3: Collection selection versus evidence retrieval}
\label{sec:routing-audit}
In the Main comparison, the scripted selector chose the labelled collection for
all 72 in-domain questions with full route cards (A0), compared with 49
with nominally name-only cards (A3). The corresponding evidence packs
contained the target string for 36 and 26 questions, respectively
(Figure~\ref{fig:routing}). The runner logged a shelf request but chose documents by vector rank, not
from shelf descriptions. A3 could fall back to descriptions without logging
that path and could search outside the benchmark libraries. Some route cards
also share title terms with the questions. A3 returned no candidate on
35 of 72 questions, including 23 of the 49 it routed correctly. A0 found
candidates on 21 of those 23. The comparison does not isolate descriptions
or shelf navigation. Earlier runs on the same questions appear separately
in Figure~\ref{fig:routing}.

\begin{figure}[!t]
\centering
\includegraphics[width=\columnwidth]{fig_latest_route}
\caption{(a) Labelled-collection selection and (b) target-string containment
for A0 and A3. Shading marks the Main comparison. All three campaigns use
the same 72 in-domain questions; they are not independent test sets.}
\label{fig:routing}
\end{figure}

%% Evaluation results for the experience store and the content-layer arms.
\subsection{E4: Operational knowledge and automated distillation}
\label{sec:e3}
HMRS stores operational lessons as searchable documents and ranks them by
credibility (Section~\ref{sec:org-experience}). The API contract
check recorded creation of a lesson through the summarisation endpoint.
This checks the write path, not the quality or later use of the lesson.
Historical smoke tests and their limits are documented in Table~S9.

\subsection{E5: Content-layer configuration comparison}
\label{sec:ladder-opt}
A0 passes candidate references to a judge with a nominal
40{,}000-character document window. A0f explicitly reads up to 250{,}000
characters per candidate and passes the text inline. A0h uses the A0f
route, candidate search and inline judging, then runs an extra vector
search. It gives up to eight distinct chunk windows priority in a
60{,}000-character pack and fills the remaining space with judged text.
PD searches all ten libraries separately and uses inline judging. A1b
merges the ten scoped vector lists and uses reference-mode judging; A2
concatenates snippets from the same union without judging. A3 uses the
nominally name-only route cards. Routed document search requests 30 hits at
a cosine threshold of 0.35 and retries at 0.30 if empty; the cross-library
arms do not use this retry. Each arm takes up to six distinct documents
from its candidate list. The reference-mode arms can use the gold domain
to resolve candidate identifiers (Section~\ref{sec:threats}).

Table~S12 reports the seven arms and Table~S13 gives paired contrasts.
In the Main comparison, A0h matched 48 of 72 targets (66.7\%; Wilson 95\%
interval 55.2--76.5\%), versus 36 for A0 (50.0\%): 12 A0h-only hits and no
A0-only hits (two-sided exact McNemar $p=0.00049$, unadjusted). Against A0f,
the closer inline-judging comparison, A0h matched seven additional targets
without losing any (48 vs.\ 41; $p=0.01563$). This contrast tests the extra
chunk search, chunk priority and pack cap together. PD reached 42 of 72
(seven A0h-only versus one PD-only, $p=0.0703$), which does not establish
superiority. Figure~\ref{fig:content} shows the other campaigns separately:
they used the same questions but not a frozen corpus and index across runs.
The A0h-over-A0 paired direction held in all three campaigns, with 15, 12
and 12 A0h-only questions and no A0-only question in any record.
The attempted A3 rebuild stalled. The retry reused the description-stripped
document libraries but created a new route-card library.

A0h and A0 recorded mean retrieval tool times of 15.0 and 16.6~s, excluding
answer generation. The trace recorded 298.9 and 33.4~kB per query, but the
counts cover different stages: A0 includes candidate snippets and the
returned pack, omitting internal judge reads. A0h counts explicit reads,
the judged pack and chunk text before the final cap. Neither number
measures total I/O, model-token use or final context size. The historical
read-all scan found 72/72 targets at 24.3~MB per question in the Frozen
aggregate. It was incomplete in the Follow-up and omitted from the Main
comparison.

In the Main comparison, A0, A0f and A0h all routed the three \texttt{F-hlm} questions
correctly but recorded no top candidate, retained evidence or target hit.
This distinguishes an empty retrieval from a routing error; the trace does
not diagnose its cause. Across the seven arms, none of the 28 arm--probe
rows met the runner's empty-candidate or empty-survivor rule. No answer was
generated, so these probes do not measure the ability to abstain.

Figure~\ref{fig:case-anatomy} shows two cases from a separate response-capture
run, not the texts underlying the Frozen aggregate.
For the five case-arm paths displayed, the capture's containment flags agree
with the frozen trace; this agreement is not a claim about other questions.
For the question on RoBERTa, the captured A0f and A0h packs both contain
the same mention in Lewis et~al.'s RAG paper, a different document from
the annotated source. This case does not require the added chunks. For the \emph{Pride and Prejudice} question, captured A0's top
candidate is part~2, but its sole survivor is part~6; captured A0f's survivors are
parts~21 and~6. Neither captured response contains the exact target string. A0h
includes the requested opening dialogue in part~2, but a line break splits
the target phrase there. Its contiguous exact match occurs instead in a
later dialogue in part~24. Part~24 therefore cannot support the specific
claim that those words occur in the opening exchange. The opening location
is part~2; inspecting that text resolves the example despite the line break.
This illustrates how to check a passage against a claim, not answer accuracy.

\begin{figure*}[!b]
\centering
\includegraphics[width=\textwidth]{fig_opt_case}
\caption{Evidence paths from a separate response-capture run. The five displayed case-arm
containment flags agree with the Frozen aggregate. A0h combines inline judged
text with up to eight chunk windows, taking chunks first under a
60{,}000-character cap; this is a configuration comparison.
(a) A0h includes a passage
containing \emph{RoBERTa} in the paper by Lewis et~al.~\citep{rag-lewis}.
The bracketed reference in the excerpt uses the source paper's numbering;
the captured A0f pack (not drawn) contains the same mention.
(b) A0h's exact match for \emph{four or five thousand a year} occurs in
later dialogue (part~24); its opening dialogue (part~2) splits the phrase
across a line break. HIT and MISS refer to exact target-string containment
in the captured pack, not answer correctness.}
\label{fig:case-anatomy}
\end{figure*}

\begin{figure*}[!t]
\centering
\includegraphics[width=\textwidth]{fig_latest_content}
\caption{Target-string containment on the same 72 in-domain questions.
(a) Filled circles and Wilson 95\% intervals show the Main comparison;
open squares and triangles show the Frozen aggregate and Follow-up.
(b) Paired outcomes in the Main comparison: A0-only hits are on the left
and comparison-arm-only hits are on the right. Questions with the same
outcome do not contribute to the bars. Read-all was not part of this
comparison. Containment does not establish answer correctness.}
\label{fig:content}
\end{figure*}

%% Answer diagnostics, deployment acceptance and historical boundary cases.
\subsection{E6: Answer diagnostics}
\label{sec:e5}
\label{sec:agentic}\label{sec:abstention}\label{sec:composition}\label{sec:eval-agent}
In the platform-only check, the system answered ten paper-corpus questions
with the five relevant collections predefined. All ten responses completed
and included the annotated paper identifier; nine contained all designated
terms. Mean wall time was 131.2~s and total recorded cost was \$3.10. The
broad marker rule counted all ten responses, but it also recognises bare
identifiers. Masking the gold identifier left a marker in eight answers;
a separate rule requiring an explicit repository label or link found three.
Neither rule tests whether a citation resolves or supports a claim. This
check used the original BQ06 prompt and had no concurrent file or chunk
comparison. Tables~S14--S15 report the original and common-library grids
as separate diagnostics.

\subsection{E7: Deployment checks}
\label{sec:sysarms}
A one-document API fixture recorded six passing assertions.
The keyword check matched the document name and the catalogue check found
its collection. The vector assertion could pass on a response with an empty
results list, so the pass flag does not establish a vector hit. The contract
suite initially passed 11 of 13 checks; the two document-read checks passed
when repeated after fixture creation. Separate answer probes used direct
reads when the verifier was unavailable; they do not retest the verifier.
These checks establish selected interface behaviour, not retrieval accuracy.
In the separate catalogue-truncation case, an incomplete preview blocked
an unpinned search and pagination recovered the affected query.

\subsection{E8: Historical candidate-ordering diagnostic}
\label{sec:e6}
\label{sec:inflation}\label{sec:eval-verify}
The archived industrial-corpus diagnostic illustrates why admission and
ordering should be evaluated separately: sorting retained candidates by
verifier score moved some gold documents outside the read window.
Figures~S4--S5 and S7 report the historical ranks, separate score archive
and limits of the replay. The deployment fixtures do not repeat this
comparison or estimate the effect of the admission threshold.

\subsection{E9: Operational boundaries}
\label{sec:e7}
In an earlier cross-library failure audit, global search skipped eight
unreadable index collections but reported success. Reindexing restored a
missing global hit in one tested library. A separate read-only inspection
checked for persisted records in the application and storage layers,
including session data, vector records, and collection metadata. Eleven
assertions passed. The inspection did not include a restart or recovery test
and did not retest the missed-library case.
Historical load sweeps and fixed-candidate and library-scope probes remain
separate (Table~S10 and the companion inventory). In particular, a request
unfinished at the end of a load-test observation window is not necessarily
an explicit timeout. The present evidence does not establish a current
capacity limit, concurrent-write safety or crash recovery.

\section{Discussion and lessons learned}
\label{sec:discussion}

The lifecycle audits show what a collection-level representation buys in
practice. A rename updated document state while leaving nine stale chunks
in the vector index. Re-indexing removed the duplicates and restored the
expected count. The three-source checks let operators verify the file,
registry and metadata views separately, while a search probe checks whether
the index reflects the change. The paper-ingestion audit reports separate
observations for a derived library and the five-library inventory
(Section~\ref{sec:e1}). The practical implication is to
test retrieval after a lifecycle operation as well as checking its write
response. This joins metadata governance to a property that matters at
query time \citep{park-metadata-quality-2009,paivarinta-ecm-2005}.

At the access layer, the full route cards gave the scripted selector a
usable collection-level view. The evidence packs show a second decision
after routing: how much of the selected material to read. Adding document
and vector chunk windows in A0h raised target-string containment over A0 by
15 questions in the frozen run and 12 in each same-question rerun. The
frozen run used an average of 309~kB of trace-recorded text per question,
compared with 31.4~kB for A0 and 24.3~MB for the mechanical full read. Thus
A0h traded a larger reading budget for more retained targets while using
about 1.3\% of the full scan's text volume. A0h also led the per-library
union PD by six questions in both the frozen and November~12 records, with a smaller paired difference
that remained unresolved. PD searches each of the ten libraries separately
before merging document candidates. The stronger result is the improvement over A0
when the retrieval configuration expands the evidence it reads.

Source visibility carries through to answering. On different substrates,
the rates of detected document-style markers diverged across the three
interfaces. Counts were nine of ten for each arm in the repeat using the same
indexed libraries, but the changed BQ06 wording and access configurations
prevent attribution of that convergence to the substrate alone. Marker checks are a
trace-format diagnostic. A separate semantic assessment would be needed to
determine which answer claims the retrieved passages support.

The historical ordering comparison identifies a distinct design choice.
Among retained candidates, fusion order produced 15 top-five hits across
sixteen queries, versus seven in verifier-score order. The separate score
archive contains only scores above the threshold; it does not establish the
effect of active filtering. This distinction parallels the practical
separation of relevance judgement and ranking in work on
LLM judgements and reranking \citep{zheng-llmjudge,sun-urara}.

Operational traces also show why a search needs to report its coverage.
Pagination recovered a failed catalogue path. An index read failure exposed
a different defect: the global channel could skip a library and still report
success. For the affected example, re-indexing recovered the missing hit.
Reporting partial coverage would help callers distinguish an incomplete search
from a completed search that returned no supporting evidence.

\section{Threats to validity}
\label{sec:threats}

\noindent\textbf{Data and scope.} The collections represent local
demonstration and benchmark workloads. Developers authored or expanded the
questions and gold annotations. The routing benchmark includes 52 expanded
questions, with 47 questions concentrated in three English domains.
Same-question reruns test repeatability under recorded changes to the
environment. An independent question set would test transfer. The storage
reconciliation examined 33 of 322 parts, and the four out-of-domain probes
characterise individual responses.

\noindent\textbf{Measurement and comparisons.} The content metric counts a
literal target string anywhere in the pack, including in a sibling
document. A post-hoc split by case-insensitive question--target overlap
(Section~\ref{sec:eval-setup}) gives frozen A0 and A0h containment of
16/29 and 21/29 on overlapping questions, and 21/43 and 31/43 on the
remainder. These are descriptive, non-random subsets, not a causal test of
overlap; lack of an exact match does not rule out other lexical cues.
Gold-ID checks search answers and recorded trace text; the saved grid
questions do not contain their corresponding IDs. Nevertheless, ID presence
does not establish source resolution or semantic correctness, and
claim-to-source support requires its own evaluation; citation markers only
indicate recognised formatting. The full-read comparator is a mechanical
scan of the corpus. Content arms change several components, including candidate
construction and reading scope. A0, A1b and A3 can use the gold domain when
resolving references, and A3 can fall back to descriptions or retrieve
outside the evaluated libraries. The archived traces leave the frequency of
these paths unknown. The reported contrasts therefore apply to the tested
configurations. A matched intervention would be needed to attribute
differences to navigation alone. The two ten-question answer grids use one
model and one run per arm, with the common-substrate grid controlling which
libraries each interface uses.

\noindent\textbf{Statistical and operational validity.} The paired
comparisons are exploratory. Wilson intervals summarise uncertainty
conditional on the question set and observed run. Candidate- and span-level
tests can understate uncertainty by treating within-query observations as
independent. Some historical reference-mode inputs, judge inputs and code
hashes are missing. The verifier results use a historical implementation
and support replay of the released records. The lifecycle assertions and
concurrency sweep cover one installation. The experience approval path,
multi-host reliability, and crash recovery remain open.

\section{Conclusion}
\label{sec:conclusion}

HMRS links the state of collections, documents and segments to the
interfaces used to retrieve and inspect evidence. The lifecycle
and API checks show why a successful write should be followed by a search
probe: the write response alone cannot verify index visibility, and a
cached result need not reflect current index state.

In the Main comparison, the document-plus-chunk configuration matched the
annotated string on 48 of 72 questions, versus 36 for the reference-mode
baseline and 41 for the wider document-window configuration. These are
configuration-level results on one shared question set. Correct routing did
not guarantee passage access, and a literal match did not establish that a
passage supported the answer. The results support checking collection
selection, passage access and answer support separately.

%% n_back.tex — back matter, included by main_20p.tex after n11_conclusion.tex
%% and before the bibliography. This file is the single source of truth for
%% every declaration; n11_conclusion.tex carries the Conclusion and nothing else.
%%
%% The declarations follow the Conclusion. Keep the generative-AI statement
%% immediately before the references per the current KBS Guide for Authors.
%%
%% AUTHOR ACTION: replace the CRediT names, confirm the funding wording, and
%% adjust or remove the acknowledgements.

\section*{CRediT authorship contribution statement}
\textbf{Author 1:} Conceptualisation, Methodology, Software, Investigation,
Data curation, Writing -- original draft. \textbf{Author 2:} Validation,
Writing -- review \& editing, Supervision. %% AUTHOR ACTION: adjust.

\section*{Declaration of competing interest}
The authors declare that they have no known competing financial interests
or personal relationships that could have appeared to influence the work
reported in this paper.

\section*{Funding}
This research did not receive any specific grant from funding agencies in
the public, commercial, or not-for-profit sectors. %% AUTHOR ACTION

\section*{Ethics statement}
This article does not report studies involving human participants, human
data, or animals.

\section*{Data availability}
The \href{https://github.com/kingdol666/rag-knowledge}{project repository}
provides code, benchmark scripts and a synthetic industrial-corpus recipe.
The supplement and companion inventory document results and protocols. The
trace CSVs support recomputation from saved labels, including all 532 rows
from the Main comparison, but omit question text, targets and returned
passages. The platform-only answers can be rescored locally; only derived
tables are released. Service assertions are sanitised. Missing complete
packs and index snapshots prevent live replay. Source rights limit corpus
redistribution.
%% AUTHOR ACTION: add the actual dataset repository link after deposit, or
%% explain at submission which research data cannot be shared (KBS Option C).

\section*{Acknowledgments}
The authors thank the maintainers of the open-source components the platform
builds on. %% AUTHOR ACTION: adjust or remove.

%% AUTHOR ACTION: confirm the tool used for each figure concept prototype,
%% its role, the caption disclosures and author review before submission.
\section*{Declaration of generative AI and AI-assisted technologies in the manuscript preparation process}
Anthropic Claude Code, ZCode and OpenAI Codex assisted with language editing,
manuscript revision and preparation of tables from archived results. An AI
image tool assisted with initial figure concepts. The research uses of
language models are described with the retrieval workflow and experimental
design. The authors reviewed the manuscript and take responsibility for its
content.



{\ifdefined\HMRSReviewLayout\else\footnotesize\fi
\ifdefined\HMRSReviewLayout
  \setlength{\bibsep}{0.5pt plus 0.2ex}
\else
  \setlength{\bibsep}{0pt}
\fi
\urlstyle{same}
\bibliographystyle{elsarticle-num}
\bibliography{refs,refs_new_verified}
}

\end{document}
